\documentclass[11pt,a4paper]{article}
\usepackage[margin=25mm]{geometry}
\usepackage{amsmath,amssymb,amsthm}
\setlength{\emergencystretch}{2em}
\newtheorem{lemma}{Lemma}
\newtheorem{theorem}[lemma]{Theorem}
\newcommand{\gen}[1]{\langle#1\rangle}
\newcommand{\Alt}{\operatorname{Alt}}
\title{Critical-tail transitivity in the exceptional paired branch}
\author{Human Compression project: L044 continuation}
\date{}
\begin{document}\maketitle
Retain the full critical six-point-head hypotheses, including actual
ambient IP, all maximal omissions intransitive and the smaller-degree
alternating induction. In the contrary rank range the accepted
critical-head theorem gives $n=2r\ge18$ and
\[
\begin{split}
 D&=\{1,\ldots,6\},\quad E=\Omega\setminus D,\\
 \tau&=\rho_3=(6\ b)\delta,\quad
 \sigma=\rho_4,\quad H=G_{\ge5},\quad K=G_{>5},\\
 R&=\gen{\tau,\sigma,H}=\Alt(E\cup\{6\}),\qquad
 B=\gen{\sigma,H}.
\end{split}
\]
The head on $D$ is actual and supported; $\sigma$ fixes $D$;
$H$ fixes $D\cup\{b\}$ and centralizes $\delta$.
The odd involution $\delta$ fixes $D\cup\{b\}$.
No transitivity of $B$ is assumed. All eligible edges refer to
full ambient omissions. Every predecessor proof and complete
required finite input remains an acquired premise.

\begin{lemma}[The invariant-set replacement]\label{bridge}
The point $x=b^\sigma$ differs from $b$, is fixed by $\tau$, and
is fixed by $K$. If a nonempty proper $B$-invariant subset
$S\subset E$ is also $\delta$-invariant, the configuration is
impossible.
\end{lemma}
\begin{proof}
If $\sigma$ fixed $b$, then $R$ would preserve $\{6,b\}$,
contrary to its transitivity. The full actual pure-overlap/IP
lemma gives $\tau x=x$, using the natural $R$; its proof does
not assume transitivity of $B$. Finally $K$ fixes $b$ and
centralizes $\sigma$, so fixes $x$.

If $b\notin S$, the groups $B$ and $\gen\tau$ both preserve
$S$, a proper nonempty subset of $E\cup\{6\}$.
If $b\in S$, they both preserve $S\cup\{6\}$ instead.
Either contradicts transitivity of $R$. This is not a claim
that arbitrary $B$-invariant sets are $\delta$-invariant.
\end{proof}

\begin{theorem}[Pairs assigned to actual orbits]\label{assigned}
Suppose every colour $j\ge5$ has two ambient-eligible edges in
some single actual $H$-orbit. Then the contrary rank range is
impossible, without the earlier additional assumption that $B$
is transitive on $E$.
\end{theorem}
\begin{proof}
We use the whole accepted orbit-transport, equivariant aggregation,
strict internal and localized internal proofs. Their combinatorial
and group hypotheses are unchanged: $G$ is transitive, the specified
seven-point head is fixed by $H$, $\tau$ centralizes $H$ and
belongs to each relevant omission, $n=2r$, and full actual IP and
the natural critical $R$ hold. Assign each colour exactly once
to an orbit with its pair and use the accepted normalization.
For assigned orbits $O_s$, put $d_s=|F_s|-|O_s|/2$, let $f$
count $H$-fixed points and $c$ its nontrivial orbits. The complete
paid calculation gives
\[
 \sum_s d_s=f/2-5\ge-3/2,\qquad f+c\le10.
\]
It excludes odd exchanged nontrivial orbits, nonfirst internal
odd orbits, and empty internal odd orbits of size at least six.
The remaining nonempty first-assigned budgets are precisely
$(f,c)=(8,1)$ or $(7,2)$, with an internal orbit of size $2q$,
$q$ odd; the latter also has the odd anchored orbit and final
localized interval of the accepted proof.

Here are all changed transitivity interfaces in that argument.
In the $(8,1)$ case let $z$ be the sole additional $H$-fixed point
outside $D\cup\{b\}$. The involution $\delta$ fixes $z$ because
it centralizes $H$ and preserves the other seven fixed points.
If $x=z$, then $S=\{b,z\}$ is preserved by $B$ and $\delta$,
contradicting Lemma~\ref{bridge}. Otherwise $x$ lies in the free
internal orbit, contrary to $\tau x=x$. In the $(7,2)$ case,
the original proof merely needs $x\ne b$, $\tau x=x$ and
$Kx=x$, all supplied by Lemma~\ref{bridge}. It places $x$ in
the odd orbit fixed pointwise by $\tau$. The complete localized
interval inequality then excludes that budget unchanged.

The remaining empty internal pair or pair of exchanged fixed points
reduces, by the paid fixed-point count, to
$\tau=(6\ b)(u\ v)$, with $H$ preserving $\{u,v\}$ and
transitive on its complementary nontrivial orbit. The two-point
rotation/localization proof uses $B$-transitivity only to obtain
$x\ne b$ and to show that $\sigma$ cannot preserve $\{u,v\}$.
The first follows from Lemma~\ref{bridge}; for the second use the
same lemma with $S=\{u,v\}$, which $\delta$ preserves.
Every later step is unchanged: moving just one point produces
incompatible simultaneous three-cycles in the actual dihedral
intersection; moving both confines all eligible later edges to
a set $Y$ with $2|I|=|Y|$, contradicting
$2|I|\le |Y|-o_{G_I}(Y)$. Singleton orbits are still counted.

This checks every use of $B$-transitivity in the acquired internal
case proofs. All possible odd contributions to the action of
$\tau$ outside its head are therefore excluded. Its exterior
action is odd, the required contradiction.
\end{proof}

\begin{theorem}[A later unique ambient colour always exists]
In every critical six-point-head configuration, some colour
$j\ge5$ has exactly one ambient-eligible edge.
\end{theorem}
\begin{proof}
Assume each such colour has at least two eligible edges. The complete
global-pair proof uses the actual orbit bound $f+c\le10$, $f\ge7$,
and the odd central action of $\tau$. We retain that full proof and
its selected-localization and maximal-string inputs. We check the
only points where its previous statement assumed $B$-transitivity.

If $c=1$, or if free-orbit equivariance puts each colour's two
edges in a single orbit, apply Theorem~\ref{assigned}.
If $c=2$ and $\tau$ exchanges the two nontrivial odd orbits,
then $f=8$ and $x$ is the additional fixed point, by
$\tau x=x$. Thus $S=\{b,x\}$ is $B$-invariant and fixed by
$\delta$, contradicting Lemma~\ref{bridge}.
This is the only changed step in that two-orbit case.

For $c=3,f=7$, the preserved free-orbit alternative uses only
the exact partition of labels and two inequalities
$|L|\le q-1$, $|F|\le(v-3)/2$, against
$r-5=q+(v-3)/2$; it is unchanged.
In the exchanged-row alternative, the complete equality argument
gives two odd rows $O,O'$ exchanged by $\tau$ and an odd orbit
$V$ fixed by $\tau$. It proves that $G_F$ is actually supported
and transitive on $V$, and that the row path labels form a
consecutive interval; no projected IP is inferred.
If the first tail label $5$ lies among the row labels, then
$\sigma$ preserves $V$. The set $S=V$ is preserved by $B$
and fixed by $\delta$, so Lemma~\ref{bridge} is the contradiction.
If $5\in F$, the same paid orbit/path argument forces the row
labels to be terminal and makes $\sigma$ preserve $O\cup O'$.
This time $S=O\cup O'$ is $B$-invariant and $\delta$ exchanges
its rows. Lemma~\ref{bridge} again applies.

The acquired case partition exhausts the possible odd exterior
actions. Hence not every colour has two eligible edges. Every
colour has at least one, since its omission is intransitive and
the whole group is transitive. The claimed unique edge follows.
\end{proof}

\paragraph{Exact reuse and source scope.}
This is an additive hypothesis extension of the exceptional
parts of \emph{Completion of the ambient paired-head branch}
and \emph{Global eligible pairs behind a six-point head}.
Their complete proofs remain paid, together with
\emph{Strict row exits and critical internal orbit budgets} and
\emph{Localized eligible edges and the final internal orbit cases}.
The displayed invariant-set replacements identify each changed
inference; none asserts that the old stronger hypothesis follows.
The source route is the later-crossing existence step in
Propositions~5.31--5.32, physical pages 38--45.
The argument does not identify the resulting cut as unshared,
settle the remaining opposite ends, or recover all of L044.
No new computation, formalization, original or question credit
is claimed.
\end{document}
